
\documentclass[a4paper,12pt]{article}

\usepackage[utf8]{inputenc}
\usepackage[T1]{fontenc}
\usepackage{lmodern,microtype}
\setlength{\emergencystretch}{2em}
\usepackage[spanish,provide=*]{babel}
\usepackage{amsmath,amssymb,amsthm}
\usepackage{graphicx}
\usepackage{xcolor}
\usepackage{enumitem,needspace}
\usepackage[top=2.5cm,bottom=2.5cm,
left=2.5cm,right=2.5cm]{geometry}

\setlength{\parindent}{0pt}
\setlength{\parskip}{7pt}
\allowdisplaybreaks

\definecolor{solutioncolor}{RGB}{15,75,110}

\newcommand{\solucion}{
\par\medskip
{\color{solutioncolor}\textbf{Solución.}}
\par
}

\newcommand{\respuesta}[1]{
\[
\boxed{#1}
\]
}

\begin{document}

\begin{flushright}
\IfFileExists{logo-espol.png}
{\includegraphics[scale=0.8]{logo-espol.png}}{}
\end{flushright}

\begin{center}
{\Large\bfseries
MÉTODOS MATEMÁTICOS PARA LAS CIENCIAS}

\vspace{0.3cm}

{\large\bfseries
Lección 1 -- Solucionario}

\vspace{0.2cm}
\end{center}

\hrule
\vspace{0.5cm}

\begin{enumerate}[leftmargin=*,
itemsep=0.8cm]

%================================================
\item \textbf{El péndulo simple
(30 puntos).}

El movimiento de un péndulo con
fricción está descrito por

\begin{equation}
m\frac{d^2\theta}{dt^2}
+\mu\frac{d\theta}{dt}
+\kappa\sin\theta=0,
\label{eq:pendulo}
\end{equation}

donde $m>0$, $\mu\geq0$ y
$\kappa>0$.

\begin{enumerate}[label=\alph*),
leftmargin=*,itemsep=0.5cm]

%------------------------------------------------
\item \textbf{(10 puntos)}
Transforme la ecuación en un
sistema de primer orden.

\solucion

Introducimos la nueva variable $v(t)=\frac{d\theta}{dt}$. Por definición, tenemos $\frac{d\theta}{dt}=v$. Derivando la expresión de $v$
respecto al tiempo, obtenemos

\[
\frac{dv}{dt}
=
\frac{d^2\theta}{dt^2}.
\]

Sustituyendo estas expresiones
en la ecuación original,

\[
m\frac{dv}{dt}
+\mu v+\kappa\sin\theta=0.
\]

Despejando la derivada de $v$,

\[
\frac{dv}{dt}
=
-\frac{\mu}{m}v
-\frac{\kappa}{m}\sin\theta.
\]

Por tanto, el sistema equivalente
de primer orden es

\[
\boxed{
\begin{aligned}
\dot{\theta}&=v,\\
\dot{v}&=
-\frac{\mu}{m}v
-\frac{\kappa}{m}\sin\theta.
\end{aligned}
}
\]

El estado del sistema queda
determinado por el vector

\[
\mathbf X(t)=
\begin{pmatrix}
\theta(t)\\v(t)
\end{pmatrix}.
\]

En particular, el sistema puede
escribirse como

\[
\dot{\mathbf X}=F(\mathbf X),
\]

donde

\[
F(\theta,v)=
\begin{pmatrix}
v\\[2mm]
-\dfrac{\mu}{m}v
-\dfrac{\kappa}{m}\sin\theta
\end{pmatrix}.
\]

Así, una ecuación diferencial
de segundo orden se transforma
en un sistema de dos ecuaciones
diferenciales de primer orden.

%------------------------------------------------
\item \textbf{(10 puntos)}
Determine todos los puntos
de equilibrio.

\solucion

Un punto de equilibrio
$(\theta_*,v_*)$ satisface $\dot{\theta}=0$ y $\dot v=0$. Utilizando el sistema obtenido
en el literal anterior,

\[
\begin{cases}
v_*=0,\\[2mm]
-\dfrac{\mu}{m}v_*
-\dfrac{\kappa}{m}
\sin\theta_*=0.
\end{cases}
\]

La primera ecuación implica $v_*=0$. Sustituyendo en la segunda,
\[
-\frac{\kappa}{m}
\sin\theta_*=0.
\]

Como $m>0$ y $\kappa>0$, $\sin\theta_*=0$. Recordemos que

\[
\sin\theta=0
\quad\Longleftrightarrow\quad
\theta=k\pi,
\qquad k\in\mathbb Z.
\]

Por consiguiente, todos
los puntos de equilibrio son

\[
\boxed{
(\theta_*,v_*)=(k\pi,0),
\qquad k\in\mathbb Z.
}
\]

\needspace{5\baselineskip}
\textbf{Interpretación física.}

Existen dos tipos de posiciones
de equilibrio:

\begin{itemize}

\item Cuando $\theta=2k\pi$,
el péndulo se encuentra en
la posición vertical inferior.

\item Cuando
$\theta=(2k+1)\pi$,
el péndulo se encuentra en
la posición vertical superior.

\end{itemize}

En ambos casos, la velocidad
angular es nula. La existencia de un equilibrio
no implica necesariamente que
este sea estable.

%------------------------------------------------
\item \textbf{(10 puntos)}
Considere $\mu=0$ y demuestre
que la energía se conserva.

\solucion

La energía mecánica total es

\[
E(\theta,v)
=
\frac12 mv^2+
\kappa(1-\cos\theta).
\]

Como $\theta$ y $v$ dependen
del tiempo, aplicamos la
regla de la cadena:

\[
\frac{dE}{dt}
=
\frac{\partial E}
{\partial\theta}
\frac{d\theta}{dt}
+
\frac{\partial E}
{\partial v}
\frac{dv}{dt}.
\]

Calculamos las derivadas
parciales:

\[
\frac{\partial E}
{\partial\theta}
=
\kappa\sin\theta,
\]

y

\[
\frac{\partial E}
{\partial v}=mv.
\]

Por tanto,

\[
\frac{dE}{dt}
=
\kappa\sin\theta\,\dot\theta
+mv\dot v.
\]

Utilizando $\dot\theta=v$,

\[
\frac{dE}{dt}
=
\kappa v\sin\theta
+mv\dot v.
\]

Factorizando,

\[
\frac{dE}{dt}
=
v(\kappa\sin\theta+m\dot v).
\]

Para el sistema original,

\[
m\dot v=-\mu v-\kappa\sin\theta.
\]

Sustituyendo,

\begin{align*}
\frac{dE}{dt}
&=
v(\kappa\sin\theta
-\mu v-\kappa\sin\theta)\\
&=-\mu v^2.
\end{align*}

Hemos demostrado la identidad
general

\[
\boxed{
\frac{dE}{dt}=-\mu v^2.
}
\]

En particular, cuando
$\mu=0$,

\[
\boxed{
\frac{dE}{dt}=0.
}
\]

Esto implica que

\[
E(\theta(t),v(t))=E_0,
\]

donde $E_0$ es una constante
determinada por las condiciones
iniciales:

\[
E_0=
\frac12 mv(0)^2+
\kappa(1-\cos\theta(0)).
\]

\needspace{5\baselineskip}
\textbf{Interpretación física.}

La energía total es la suma
de la energía cinética,

\[
K=\frac12 mv^2,
\]

y la energía potencial,

\[
U=\kappa(1-\cos\theta).
\]

En ausencia de fricción,
la energía puede transformarse
de cinética a potencial y
viceversa, pero su valor total
permanece constante.

Cuando $\mu>0$, en cambio,

\[
\frac{dE}{dt}=-\mu v^2\leq0.
\]

En este caso la energía nunca
aumenta, pues existe disipación
debida a la fricción.

\end{enumerate}

%================================================
\item \textbf{Análisis de un sistema
dinámico (35 puntos).}

Considere

\begin{equation}
\begin{cases}
\dot x=y,\\
\dot y=-4x-4y.
\end{cases}
\label{eq:sistema}
\end{equation}

\begin{enumerate}[label=\alph*),
leftmargin=*,itemsep=0.5cm]

%------------------------------------------------
\item \textbf{(10 puntos)}
Determine los puntos de equilibrio.

\solucion

Un punto de equilibrio $(x_*,y_*)$ debe satisfacer $\dot x=0$ y $\dot y=0$. Por tanto,
\[
\begin{cases}
y_*=0,\\
-4x_*-4y_*=0.
\end{cases}
\]

De la primera ecuación, $y_*=0$. Sustituyendo en la segunda, $-4x_*=0$. Por consiguiente, $x_*=0$. El sistema posee un único
punto de equilibrio:
\[
\boxed{(x_*,y_*)=(0,0).}
\]

%------------------------------------------------
\item \textbf{(10 puntos)}
Escriba el sistema en
forma matricial.

\solucion

Definimos el vector de estado

\[
\mathbf X=
\begin{pmatrix}
x\\y
\end{pmatrix}.
\]

Su derivada temporal es

\[
\dot{\mathbf X}=
\begin{pmatrix}
\dot x\\\dot y
\end{pmatrix}.
\]

Utilizando las ecuaciones
del sistema,

\[
\dot{\mathbf X}=
\begin{pmatrix}
y\\-4x-4y
\end{pmatrix}.
\]

Observamos que

\[
\begin{pmatrix}
y\\-4x-4y
\end{pmatrix}
=
\begin{pmatrix}
0&1\\
-4&-4
\end{pmatrix}
\begin{pmatrix}
x\\y
\end{pmatrix}.
\]

Por tanto,

\[
\boxed{
\dot{\mathbf X}
=A\mathbf X,
\qquad
A=
\begin{pmatrix}
0&1\\
-4&-4
\end{pmatrix}.
}
\]

Esta representación permite
estudiar la dinámica del
sistema mediante las propiedades
de la matriz $A$.

%------------------------------------------------
\item \textbf{(10 puntos)}
Calcule los valores propios
y clasifique el equilibrio.

\solucion

Los valores propios de $A$
son las soluciones de la
ecuación característica

\[
\det(A-\lambda I)=0.
\]

Primero calculamos

\[
A-\lambda I=
\begin{pmatrix}
-\lambda&1\\
-4&-4-\lambda
\end{pmatrix}.
\]

Su determinante es

\begin{align*}
\det(A-\lambda I)
&=(-\lambda)(-4-\lambda)
 -(1)(-4)\\
&=\lambda(4+\lambda)+4\\
&=\lambda^2+4\lambda+4.
\end{align*}

Factorizando,

\[
\lambda^2+4\lambda+4
=(\lambda+2)^2.
\]

Entonces,

\[
(\lambda+2)^2=0.
\]

Se obtiene un valor propio
real repetido:

\[
\boxed{\lambda_1=\lambda_2=-2.}
\]

Para clasificar correctamente
el equilibrio, verificamos
los vectores propios.

Resolvemos

\[
(A+2I)\mathbf u=0.
\]

Como

\[
A+2I=
\begin{pmatrix}
2&1\\
-4&-2
\end{pmatrix},
\]

tenemos

\[
\begin{pmatrix}
2&1\\
-4&-2
\end{pmatrix}
\begin{pmatrix}
u_1\\u_2
\end{pmatrix}
=
\begin{pmatrix}
0\\0
\end{pmatrix}.
\]

Este sistema se reduce a

\[
2u_1+u_2=0.
\]

De aquí,

\[
u_2=-2u_1.
\]

En consecuencia, los vectores
propios tienen la forma

\[
\mathbf u=s
\begin{pmatrix}
1\\-2
\end{pmatrix},
\qquad s\neq0.
\]

Existe solamente una dirección
propia linealmente independiente.

Por consiguiente, la matriz
no es diagonalizable.

Como el valor propio repetido
es negativo, el equilibrio
es un

\[
\boxed{
\text{nodo estable degenerado}
}
\]

también llamado
\textbf{nodo impropio estable}.

En particular, el origen
es asintóticamente estable.

%------------------------------------------------
\item \textbf{(5 puntos)}
Describa el comportamiento
de las trayectorias.

\solucion

Como ambos valores propios
tienen parte real negativa,

\[
\lambda_1=\lambda_2=-2<0,
\]

las soluciones que comienzan
cerca del origen se aproximan
al equilibrio cuando
$t\to\infty$.

Es decir,

\[
\boxed{
\lim_{t\to\infty}
\mathbf X(t)
=
\begin{pmatrix}
0\\0
\end{pmatrix}.
}
\]

En consecuencia,

\[
\lim_{t\to\infty}x(t)=0,
\qquad
\lim_{t\to\infty}y(t)=0.
\]

De hecho, por tratarse de
un sistema lineal con todos
sus valores propios de parte
real negativa, la convergencia
ocurre para cualquier condición
inicial.

\textbf{Verificación adicional.}

La primera ecuación implica

\[
y=\dot x.
\]

Derivando,

\[
\dot y=\ddot x.
\]

Sustituyendo en la segunda
ecuación,

\[
\ddot x+4\dot x+4x=0.
\]

La ecuación característica es

\[
r^2+4r+4=0,
\]

que tiene la raíz doble $r=-2$.

La solución general es

\[
x(t)=(C_1+C_2t)e^{-2t}.
\]

Derivando,

\[
\begin{aligned}
y(t)
&=\dot x(t)\\
&=\left[
C_2-2C_1-2C_2t
\right]e^{-2t}.
\end{aligned}
\]

Como

\[
e^{-2t}\to0,
\qquad
te^{-2t}\to0,
\]

obtenemos

\[
x(t)\to0,\qquad y(t)\to0.
\]

Esto confirma la estabilidad
asintótica del origen.

\end{enumerate}

%================================================
\item \textbf{Modelación de una
reacción química (35 puntos).}

Considere la reacción irreversible

\[
A\xrightarrow{k}B,
\qquad k>0.
\]

Denotamos por $a(t)$ y $b(t)$
las concentraciones de las
sustancias $A$ y $B$.

La velocidad de transformación
es proporcional a la
concentración de $A$.

\begin{enumerate}[label=\alph*),
leftmargin=*,itemsep=0.5cm]

%------------------------------------------------
\item \textbf{(10 puntos)}
Formule el sistema de ecuaciones
diferenciales.

\solucion

La hipótesis de proporcionalidad
indica que la velocidad de
reacción es

\[
r(t)=ka(t),
\]

donde $k>0$ es la constante
de reacción.

Como la sustancia $A$ se consume,
su concentración disminuye
con una velocidad igual a
$r(t)$.

Por tanto,

\[
\frac{da}{dt}=-r(t)=-ka(t).
\]

Por otro lado, la sustancia
$B$ se produce a partir de $A$.

Suponiendo una reacción
uno a uno, cada unidad de
$A$ consumida genera una
unidad de $B$.

Por tanto,

\[
\frac{db}{dt}=r(t)=ka(t).
\]

El modelo matemático es

\[
\boxed{
\begin{cases}
\dot a=-ka,\\
\dot b=ka.
\end{cases}
}
\]

Se trata de un sistema
autónomo y lineal de dos
ecuaciones diferenciales
de primer orden. La primera ecuación describe
la desaparición de $A$,
mientras que la segunda
describe la formación de $B$.

%------------------------------------------------
\item \textbf{(10 puntos)}
Demuestre que $a(t)+b(t)$
permanece constante.

\solucion

Definimos la concentración
total

\[
C(t)=a(t)+b(t).
\]

Derivando respecto al tiempo,

\[
\frac{dC}{dt}
=
\frac{da}{dt}
+
\frac{db}{dt}.
\]

Sustituyendo las ecuaciones
del modelo,

\[
\frac{dC}{dt}
=
-ka+ka = 0
\]

Es decir,

\[
\boxed{
a(t)+b(t)=a(0)+b(0).
}
\]

Esta identidad constituye
una ley de conservación
para el sistema.

\needspace{5\baselineskip}
\textbf{Interpretación física.}

La reacción transforma
la sustancia $A$ en la
sustancia $B$ sin crear
ni destruir la cantidad
total representada por
$a+b$. La cantidad de $A$ que
desaparece es exactamente
igual a la cantidad de
$B$ que se produce. Esto supone una reacción
con estequiometría uno a uno,
sin pérdidas ni fuentes
externas y con volumen
constante.

%------------------------------------------------
\item \textbf{(10 puntos)}
Resuelva el sistema con

\[
a(0)=a_0>0,\qquad b(0)=0.
\]

\solucion

Comenzamos con la primera
ecuación:

\[
\frac{da}{dt}=-ka.
\]

Es una ecuación diferencial
de primer orden con
variables separables.

Como $a_0>0$, podemos
escribir

\[
\frac{1}{a}\frac{da}{dt}=-k.
\]

Separando las variables,

\[
\frac{da}{a}=-k\,dt.
\]

Integrando,

\[
\int\frac{da}{a}
=
-k\int dt.
\]

Obtenemos

\[
\ln|a|=-kt+C.
\]

Aplicando la función
exponencial,

\[
a(t)=Ce^{-kt},
\]

donde hemos utilizado que
$a(t)>0$ para $t\geq0$. Aplicando la condición inicial,

\[
a(0)=C=a_0.
\]

Por tanto,

\[
\boxed{a(t)=a_0e^{-kt}.}
\]

Ahora resolvemos la segunda
ecuación:

\[
\frac{db}{dt}=ka(t).
\]

Sustituyendo la solución
obtenida,

\[
\frac{db}{dt}
=ka_0e^{-kt}.
\]

Integramos desde $0$ hasta $t$:

\[
b(t)-b(0)
=
\int_0^t ka_0e^{-ks}\,ds.
\]

Como $b(0)=0$,

\[
b(t)
=
ka_0\int_0^t e^{-ks}\,ds.
\]

Calculando la integral,

\[
b(t)
=
ka_0
\left[
-\frac{1}{k}e^{-ks}
\right]_0^t.
\]

Entonces,

\[
\begin{aligned}
b(t)
&=-a_0(e^{-kt}-1)\\
&=a_0(1-e^{-kt}).
\end{aligned}
\]

En consecuencia,

\[
\boxed{
b(t)=a_0(1-e^{-kt}).
}
\]

La solución completa es

\[
\boxed{
\begin{aligned}
a(t)&=a_0e^{-kt},\\
b(t)&=a_0(1-e^{-kt}).
\end{aligned}
}
\]

\textbf{Verificación.}

Sumando las soluciones,

\[
\begin{aligned}
a(t)+b(t)
&=a_0e^{-kt}
+a_0(1-e^{-kt})\\
&=a_0.
\end{aligned}
\]

Así,

\[
a(t)+b(t)=a_0,
\]

lo cual coincide con la ley
de conservación demostrada
en el literal anterior.

%------------------------------------------------
\item \textbf{(5 puntos)}
Determine los límites de
las concentraciones e
interprete los resultados.

\solucion

Recordemos las soluciones

\[
a(t)=a_0e^{-kt}
\]

y

\[
b(t)=a_0(1-e^{-kt}).
\]

Como $k>0$,

\[
\lim_{t\to\infty}e^{-kt}=0.
\]

En consecuencia,

\[
\begin{aligned}
\lim_{t\to\infty}a(t)
&=a_0\lim_{t\to\infty}e^{-kt}\\
&=0.
\end{aligned}
\]

Por tanto,

\[
\boxed{
\lim_{t\to\infty}a(t)=0.
}
\]

Para la segunda sustancia,

\[
\begin{aligned}
\lim_{t\to\infty}b(t)
&=a_0\left(
1-\lim_{t\to\infty}e^{-kt}
\right)\\
&=a_0.
\end{aligned}
\]

Por consiguiente,

\[
\boxed{
\lim_{t\to\infty}b(t)=a_0.
}
\]

\needspace{5\baselineskip}
\textbf{Interpretación física.}

A medida que transcurre
el tiempo, la sustancia
$A$ se consume gradualmente
y su concentración se
aproxima a cero. Simultáneamente, la
concentración de $B$
aumenta hasta aproximarse
a la concentración inicial
de $A$. En el límite, toda la
sustancia inicial se
encuentra transformada
en $B$. La evolución del sistema
es consistente con
la conservación de la
concentración total:

\[
a(t)+b(t)=a_0.
\]

\end{enumerate}

\end{enumerate}

%================================================


\end{document}
